\scanpage{92}

\deleted{\unsure{Il}} Montrons que a) implique c). Soit en effet \deleted{\unsure{donc}} $E$ un vectoriel \deleted{sur $Q$} \inserted{``de degré $i$''} avec $T$ opérant.

\deleted{Alors les \unsure{opér.} \unsure{de} $T_{\mathbb R}$ \unsure{opèrent} sur $E_{\mathbb R}$, et comme}
\deleted{Soit $U$ le plus \unsure{grand sous-groupe}}

Alors $T_{\mathbb R}$ opère sur $E_{\mathbb R}$. Soit $U$ le plus grand sous-tore compact de $T_{\mathbb R}$. Il résulte aisément de l'hyp. a) que $T_{\mathbb R}^{\unsure{\circ|0}}=\unsure{U_{\mathbb R}|U}\cdot S_{\mathbb R}$, où $S=\operatorname{Im}j$, \deleted{et} \inserted{et $U\subset\operatorname{Ker}\varepsilon$.} Donc $f$ est de la forme $uj(x)$, ou du moins il existe un entier \deleted{positif} \inserted{$\nu>0$ tel} que
\[
f^\nu=u\cdot j(x)\qquad u\in U(\mathbb R),\quad x\in\mathbb R^*.
\]
Or $U$ étant compact, les valeurs propres de $\forall$ $u\in U(\mathbb R)$ sont de val. absolue 1, donc celles de $f^\nu$ sont de val. absolue égale à celles de $j(x)$, or $j(x)$ est le \deleted{\unsure{représentation}} \unsure{l'hom.} \unsure{h} par $x^i$, \inserted{d'où les val. propres de $f$ sont de val. abs. $x^{i/\nu}$.} D'autre part \deleted{$\varepsilon(u)=\unsure{0}$} $\varepsilon|U=\unsure{0}$ donc
\[
\varepsilon(f)^\nu=\varepsilon(u)\varepsilon j(x)=\varepsilon j(x)=x^\alpha,
\quad\text{donc}\quad x^{i/\nu}=\varepsilon(f)^{i/\alpha}=q^{i/\alpha},
\]
d'où la conclusion.

\margin{?}
\deleted{N.B. Évidemment b) et c) n'impliquent pas a), \unsure{pour aucun $\alpha$}, car il y a des sous-groupes libres de $Q^*$ de rang $\geq2$ \ldots}

D'autre part, on sait que a équivaut au fait que $\mathcal M$ est (avec ses deux structures suppl. provenant de $j,\varepsilon$) munissable d'une structure ``positive'' du style \unsure{motivique}. Ceci, d'autre part, par un argument connu, implique que

\mspage{91}
